\section{Воспроизводимость и пределы вычислительного сертификата}
\label{sec:reproducibility}

\paragraph{Точная арифметика.}
Новый код использует Python и SymPy с рациональными коэффициентами
и явным соотношением $a^2-a+1=0$. Числа с плавающей точкой не
используются. Комментарии, сообщения об ошибках и описания функций
написаны по-английски. Проверки основаны на явных исключениях и
остаются включёнными при запуске Python с флагом \texttt{-O}.

\paragraph{Что проверяется вычислением.}
\begin{enumerate}
\item Обе конгруэнции для $P_x,F_x$, точные деления, компактная
базовая обратная карта и подстановки в трилистник.
\item Гармоническое тождество, перенесённая пара, точные деления
в прямом и обратном направлениях и восстановление кривой из оси.
\item Для явно указанного конечного диапазона индексов:
рекурсия $p_N$, полиномиальность $C_N$, обратимость центра по модулю
полного $x^N$, построение поправки, сертификаты полиномиальности
обеих карт и точные подстановки в параметризации.
\item Нормировка внутренней координаты в доказательстве дикости,
порядок её квадратичного коэффициента по $x$ и мономиальная чётность,
используемая в кремоновом доказательстве.
\item Отрицательные контрольные примеры: испорченная конгруэнция
и числитель, не допускающий требуемого точного деления, должны быть
обнаружены, включая режим \texttt{-O}.
\end{enumerate}

\paragraph{Граница вычислительного сертификата.}
Для больших карт код использует факторизованные сертификаты:
полиномиальность обеих сторон и тождества после локализации. Это
эквивалентно доказательству глобальной обратимости в области целостности;
полное раскрытие всех шести композиций в один огромный полином
не требуется и не заявляется для каждого $N$.
Утверждение для всех $N$ следует из индукции и доказательств в тексте,
а не из конечного диапазона тестов. Общие теоремы о дикости,
стабильной ручности и мономиальных преобразованиях используются
как опубликованные математические результаты. Прилагаемые программы
не являются их формальной верификацией в системе доказательств.
Явное слово в элементарных автоморфизмах шести переменных здесь
не вычислялось: существование такого слова следует из \cite{BEW2012}.

\paragraph{Запуск нового приложения.}
Из каталога \texttt{appendices/2026-10-28-cremona-program}:
\begin{verbatim}
python -m pip install -r requirements.txt
python computations/run_all.py
\end{verbatim}
Команда запускает обычный и оптимизированный режимы; точный диапазон,
версии программ, результаты и времена записываются в \texttt{results/}.
В поставляемом запуске проверены $N=1,2,3$: выполнено по 90 проверок
в каждом режиме, и их математические результаты совпали.
Описание отдельных проверок дано в \texttt{computations/README.md}.
Загружать данные из сети во время проверки не требуется.

\paragraph{Повторная проверка исходного препринта.}
Из корня полного репозитория выполнены
\begin{verbatim}
python -u verify.py
python -u verify_printed_formulas.py
\end{verbatim}
Оба запуска завершились успешно. Второй читает реальные правые части
формул исходного \LaTeX{} и проверяет все двенадцать компонент
прямых и обратных выпрямителей. Журналы находятся в
\texttt{results/upstream\_compact.log} и
\texttt{results/upstream\_printed.log} нового приложения.
Остальные старые журналы включены в полный архив как исторические
материалы; их наличие не означает, что каждый старый поиск, GAP-скрипт
или Macaulay2-скрипт был заново выполнен при подготовке Appendix 2.

\paragraph{Сборка и аудит файлов.}
Для \LaTeX{} используется \texttt{latex/main.tex}, компилятор XeLaTeX
и два прохода. Архив для Overleaf содержит все разделы, библиографию
и необходимые текстовые шрифты с лицензией. Преамбула сохраняет
формат исходного препринта: \texttt{article}, 11\,pt, A4, поля 24\,mm,
математические окружения, колонтитулы и кликабельные ссылки.
Файлы \texttt{PROVENANCE.md}, \texttt{audit/claim\_status.md} и
\texttt{SHA256SUMS.txt} описывают исходную версию, границы проверки
и контрольные суммы. Исходный основной текст и Appendix 1 не
переписываются. Это отдельная исследовательская ветка для дальнейших
исправлений и совместной работы.
